% SOURCE_AUTHORITY: sga5_fr_workpass.tex, lines 1153--1192 (LF-normalized unit).
\subsection*{4.6}
Sean $X$ un esquema, $j:Y\to X$ un subesquema cerrado e $i:\ol{x}\to Y$ un punto geométrico de $Y$. Sea, por otra parte, $K_X$ un complejo dualizante sobre $X$, normalizado en $\ol{x}$ (cf.~(4.5)). Hemos visto en (4.5) que, si $Y$ es la clausura de $x$, se tiene, para $F\in\ob\D_c^b(X)$, un apareamiento perfecto
\[
\R^{!}j(F)_{\ol{x}}\times \underline{D}_X(F)_{\ol{x}}\longrightarrow A.
\]
No es difícil definir un apareamiento análogo en el caso general. Introduzcamos la factorización de $i$ en
\[
\ol{x}\xrightarrow{i'}\ol{\{x\}}\xrightarrow{i''}Y.
\]
Para $F\in\ob\D_c^b(X)$, se tienen isomorfismos canónicos:
\[
\begin{aligned}
\R\Gamma_{\ol{x}}j^*\underline{D}_X(F)&\xrightarrow{\sim}\R\Gamma_{\ol{x}}\underline{D}_Y\R^{!}j(F)\qquad\text{(fórmula de inducción complementaria (4.2.1))}\\
&\xrightarrow{\sim}i'^*\R^{!}i''\underline{D}_Y\R^{!}j(F)\qquad\text{(definición de $\R\Gamma_{\ol{x}}$ (4.5.1))}\\
&\xrightarrow{\sim}i'^*\underline{D}_{\ol{x}}i''^*\R^{!}j(F)\qquad\text{(inducción ordinaria (1.12 b) (i))}\\
&\xrightarrow{\sim}\underline{D}_{\ol{x}}i^*\R^{!}j(F)\qquad\text{(por 4.3)), es decir, finalmente un}
\end{aligned}
\]
isomorfismo canónico
\[
(4.6.1)\qquad \R^{!}j(F)_{\ol{x}}\xrightarrow{\sim}\underline{D}_{\ol{x}}\R\Gamma_{\ol{x}}j^*\underline{D}_X(F),
\]
de donde se obtienen los apareamientos perfectos
\[
(4.6.1)'\qquad \R^{!}j(F)_{\ol{x}}\times \R\Gamma_{\ol{x}}j^*\underline{D}_X(F)\longrightarrow A,
\]
\[
(4.6.1)''\qquad \uH_Y^i(F)_{\ol{x}}\times \uH_{\ol{x}}^{-i}(\underline{D}_X(F)|Y)\longrightarrow A.
\]

\begin{statement}{Ejemplo 4.6.2}
Sean $X$ un esquema regular, $\ol{x}$ un punto geométrico de $X$, y $d=\dim(\cO_{X,\ol{x}})$. Pongamos $K_X=(\mu_n)_X^{\otimes d}[2d]$. Si $X$ satisface la condición (reg) (3.1.1) y el teorema de pureza (3.1.4), entonces $K_X$ está canónicamente normalizado en $\ol{x}$. Por tanto, si $K_X$ es dualizante, se tienen apareamientos perfectos
\[
\uH_Y^i(F)_{\ol{x}}\times \uH_{\ol{x}}^{2d-i}(F'|Y)\longrightarrow A,
\]
para todo cerrado $Y$ que contenga a $x$ y todo $F\in\ob\D_c^b(X)$, con $F'=\R\Hom(F,(\mu_n)_X^{\otimes d})$.
\end{statement}

Para terminar, demos una variante estrictamente local de (4.6.1).
